\documentclass[11pt]{article}
\usepackage[margin=1in]{geometry}
\usepackage{amsmath,amssymb,amsthm}
\usepackage{graphicx}
\usepackage{booktabs}
\usepackage{xcolor}
\usepackage{listings}
\usepackage[numbers]{natbib}
\usepackage{hyperref}

\newtheorem{theorem}{Theorem}
\newtheorem{lemma}[theorem]{Lemma}
\theoremstyle{definition}
\newtheorem{definition}[theorem]{Definition}
\theoremstyle{remark}
\newtheorem*{remark}{Remark}

\emergencystretch=2em

\lstdefinelanguage{lean}{
  morekeywords={def,theorem,have,refine,intro,induction,with,simp,only,rw,unfold,exact,example,namespace,end,open,fun,Prop,by,noncomputable,section,import},
  sensitive=true,
  morecomment=[l]{--},
  morecomment=[s]{/-}{-/},
}
\lstset{
  language=lean,
  basicstyle=\ttfamily\footnotesize,
  keywordstyle=\color{blue!60!black}\bfseries,
  commentstyle=\color{gray},
  columns=fullflexible,
  keepspaces=true,
  breaklines=true,
  frame=single,
  extendedchars=true,
  literate=
    {ℝ}{{$\mathbb{R}$}}1
    {ℕ}{{$\mathbb{N}$}}1
    {→}{{$\to$}}1
    {↔}{{$\leftrightarrow$}}1
    {∀}{{$\forall$}}1
    {∃}{{$\exists$}}1
    {∈}{{$\in$}}1
    {∑}{{$\sum$}}1
    {∏}{{$\prod$}}1
    {∧}{{$\wedge$}}1
    {≤}{{$\le$}}1
    {≠}{{$\ne$}}1
    {⨆}{{$\bigsqcup$}}1
    {ρ}{{$\rho$}}1
    {⟨}{{$\langle$}}1
    {⟩}{{$\rangle$}}1
    {·}{{$\cdot$}}1
    {𝓝}{{$\mathcal{N}$}}1
    {₀}{{$_0$}}1
    {⊢}{{$\vdash$}}1
}

\title{Loop Mass Balance for a Single-Material Recovery Loop:\\ A Lean Proof of the Fixed Point and the Virgin-Input Identity}
\author{Justin Hart, Viridis LLC}
\date{15 November 2026}

\begin{document}
\maketitle

\begin{abstract}
We study the simplest dynamic model of a circular material loop. A constant demand $D\ge0$ is
met partly by recovered material, a share $r\in[0,1]$, and partly by virgin input $V=D(1-r)$. A
secondary stock evolves as $S(t+1)=\rho S(t)+rD$ with retention $\rho\in[0,1)$. We give a Lean~4
proof, against Mathlib, of the sealed statement \texttt{P3.LoopMassBalance}. The statement has
five parts. The virgin input satisfies $0\le V\le D$. The stock gap obeys the exact closed form
$S(t)-S^\ast=\rho^t(S(0)-S^\ast)$ with fixed point $S^\ast=rD/(1-\rho)$. The stock converges to
$S^\ast$ from any initial value. Cumulative virgin input over $T$ periods equals $TD(1-r)$. For
$D>0$ and $T>0$, cumulative virgin input is strictly decreasing in $r$ on $[0,1]$. The proof is
about fifty lines long and uses only standard Mathlib lemmas. These results are accounting
identities of the model. They say nothing about any measured recovery rate, retention rate,
demand or material.
\end{abstract}

\section{Introduction}
Circular-economy policy is about closing material loops, which means substituting recovered
material for virgin extraction \citep{geissdoerfer2017}. Any quantitative tool that reasons
about such loops relies on a handful of accounting facts. Recovery reduces virgin input one for
one. A secondary stock fed at a constant rate and leaking at a constant fraction settles at a
predictable level. Raising the recovery share always lowers cumulative extraction when demand
is positive. These facts are elementary, but a decision kernel that applies them automatically
should rest on exact versions whose hypotheses are listed in full. The edge cases are where
errors hide: full retention ($\rho=1$), zero demand and out-of-range recovery fractions.

Project P3 of Viridis Research is building the formal core of such a kernel. The first session
(P3-T001, 5~November 2026) proposed four statement modules. The program owner then sealed
\texttt{P3.LoopMassBalance} under decision D-011. This note answers one question: \emph{does the
sealed loop mass balance statement hold exactly as written?} We give a Lean proof that it does.

\paragraph{Contributions.}
\begin{itemize}
\item A Lean~4 proof of \texttt{ViridisStatements.P3.LoopMassBalance.Statement}
(Section~\ref{sec:proof}). It depends only on the standard axioms \texttt{propext},
\texttt{Classical.choice} and \texttt{Quot.sound}.
\item A short human-readable proof of the same facts (Section~\ref{sec:math}). It makes clear
which hypothesis each part uses.
\item A statement of what the result settles and what it leaves open
(Section~\ref{sec:limits}).
\end{itemize}

\section{Background and related work}
\paragraph{Circular economy and industrial symbiosis.}
\citet{geissdoerfer2017} review how the circular-economy concept relates to sustainability.
Its operational core is closing material loops. \citet{chertow2000} defined industrial
symbiosis as the exchange of materials, energy, water and by-products among firms. The P3
project also has statements on cascaded recovery and $N$-partner symbiosis, which are later
tasks. The model here is the single-loop building block those statements generalise.

\paragraph{The model as a linear recurrence.}
The stock equation $S(t+1)=\rho S(t)+b$ with $b=rD$ is a first-order affine recurrence. The
fixed point, the geometric decay of the gap and convergence for $|\rho|<1$ are textbook facts.
We claim no mathematical novelty. The contribution is a machine-checkable version of exactly
the statement that the project sealed, including Lean's conventions such as $x/0=0$.

\paragraph{Formal tools.}
The proof is written in Lean~4 \citep{moura2021lean4} (toolchain v4.28.0) against Mathlib
\citep{mathlib2020} (commit \texttt{8f9d9cff}).

\section{The model and the statement}
\label{sec:model}
\begin{definition}[Loop model]
Fix demand $D$, recovery fraction $r$, retention $\rho$ and initial stock $S_0$, all real. Put
\[
V = D(1-r),\qquad S(0)=S_0,\quad S(t+1)=\rho S(t)+rD,\qquad S^\ast=\frac{rD}{1-\rho},
\qquad C(T)=\sum_{t=0}^{T-1}V .
\]
\end{definition}

\begin{theorem}[\texttt{P3.LoopMassBalance}]\label{thm:main}
\begin{enumerate}
\item If $D\ge0$, $0\le r\le1$ and $0\le\rho<1$, then for every $S_0$:
 (a) $0\le V\le D$;
 (b) $S(t)-S^\ast=\rho^t(S_0-S^\ast)$ for all $t\in\mathbb{N}$;
 (c) $S(t)\to S^\ast$ as $t\to\infty$;
 (d) $C(T)=TD(1-r)$ for all $T\in\mathbb{N}$.
\item If $D>0$ and $T>0$, then $r\mapsto C(T)$ is strictly decreasing on $[0,1]$.
\end{enumerate}
\end{theorem}

The sealed Lean text is reproduced below. In the sealed file $\mathcal{N}$ stands for
Mathlib's neighbourhood filter \texttt{nhds}.

\begin{lstlisting}[caption={Definitions and \texttt{Statement} from the sealed module \texttt{ViridisStatements.P3.LoopMassBalance} (docstrings omitted).}]
def virgin (D r : ℝ) : ℝ := D * (1 - r)

def stock (ρ r D S0 : ℝ) : ℕ → ℝ
  | 0 => S0
  | t + 1 => ρ * stock ρ r D S0 t + r * D

def steadyStock (ρ r D : ℝ) : ℝ := r * D / (1 - ρ)

def cumulativeVirgin (D r : ℝ) (T : ℕ) : ℝ := ∑ _t ∈ Finset.range T, virgin D r

def Statement : Prop :=
  (∀ D r ρ S0 : ℝ, 0 ≤ D → 0 ≤ r → r ≤ 1 → 0 ≤ ρ → ρ < 1 →
    (0 ≤ virgin D r ∧ virgin D r ≤ D) ∧
    (∀ t : ℕ, stock ρ r D S0 t - steadyStock ρ r D
        = ρ ^ t * (S0 - steadyStock ρ r D)) ∧
    Tendsto (stock ρ r D S0) atTop (𝓝 (steadyStock ρ r D)) ∧
    (∀ T : ℕ, cumulativeVirgin D r T = (T : ℝ) * D * (1 - r))) ∧
  (∀ (D : ℝ) (T : ℕ), 0 < D → 0 < T →
    StrictAntiOn (fun r => cumulativeVirgin D r T) (Set.Icc 0 1))
\end{lstlisting}

\section{Mathematical proof}
\label{sec:math}
\begin{lemma}[Fixed point and gap recursion]\label{lem:gap}
If $\rho\ne1$, then $\rho S^\ast+rD=S^\ast$. Hence $S(t+1)-S^\ast=\rho\,(S(t)-S^\ast)$.
\end{lemma}
\begin{proof}
$\rho S^\ast+rD = \frac{\rho rD + rD(1-\rho)}{1-\rho}=\frac{rD}{1-\rho}=S^\ast$. Subtracting
this from $S(t+1)=\rho S(t)+rD$ gives the recursion.
\end{proof}

\begin{proof}[Proof of Theorem~\ref{thm:main}]
(a) $V=D(1-r)$ is a product of non-negative numbers, and $D-V=Dr\ge0$.
(b) Induct on $t$ using Lemma~\ref{lem:gap}. Only $\rho<1$ is needed, through $1-\rho\ne0$.
(c) Since $0\le\rho<1$, $\rho^t\to0$. So $S(t)=S^\ast+\rho^t(S_0-S^\ast)\to S^\ast$.
(d) The sum has $T$ equal terms $V$.
(2) By (d), $C(T)$ is an affine function of $r$ with slope $-TD<0$. Note that part~(d) holds for
every real $r$, so part~2 does not need the hypotheses of part~1.
\end{proof}

\begin{remark}
The hypothesis $\rho\ge0$ is used only in (c). The convergence argument goes through for
$|\rho|<1$, but the sealed statement assumes $\rho\in[0,1)$, which is the physically meaningful
range for a retention fraction. The sealed module's negative controls show that the hypotheses
$\rho<1$, $D>0$ and $r\le1$ are needed for parts (b), 2 and (a) respectively. With $\rho=1$,
Lean's convention $x/0=0$ makes $S^\ast=0$ and the closed form fails.
\end{remark}

\section{The Lean development}
\label{sec:proof}
The candidate file \texttt{proofs/P3.LoopMassBalance.lean} follows Section~\ref{sec:math}
closely. \texttt{gap\_succ} is Lemma~\ref{lem:gap}, proved by unfolding the recurrence and
clearing the denominator $1-\rho\ne0$ with \texttt{field\_simp; ring}. \texttt{gap\_closed} is
part~(b), by induction. \texttt{stock\_tendsto} is part~(c). It combines Mathlib's
\texttt{tendsto\_pow\_atTop\_nhds\_zero\_of\_lt\_one} with \texttt{Tendsto.mul\_const} and
\texttt{Tendsto.add\_const}, then transfers along the closed form with \texttt{Tendsto.congr'}.
\texttt{cumulative\_eq} is part~(d), via \texttt{Finset.sum\_const}. The final theorem
discharges the bounds on $V$ and the strict antitonicity with \texttt{nlinarith}. For
antitonicity it is given the positivity of $TD$.

\begin{lstlisting}[caption={The complete candidate proof \texttt{proofs/P3.LoopMassBalance.lean}.}]
import Mathlib
import ViridisStatements.P3.LoopMassBalance

namespace Candidate

open Filter Topology
open ViridisStatements.P3.LoopMassBalance

/-- One-step gap recursion: `S(t+1) - S* = ρ · (S(t) - S*)`. -/
theorem gap_succ (ρ r D S0 : ℝ) (hρ : ρ < 1) (t : ℕ) :
    stock ρ r D S0 (t + 1) - steadyStock ρ r D
      = ρ * (stock ρ r D S0 t - steadyStock ρ r D) := by
  have h : (1 - ρ) ≠ 0 := by linarith
  simp only [stock, steadyStock]
  field_simp
  ring

/-- Closed form of the stock gap. -/
theorem gap_closed (ρ r D S0 : ℝ) (hρ : ρ < 1) (t : ℕ) :
    stock ρ r D S0 t - steadyStock ρ r D = ρ ^ t * (S0 - steadyStock ρ r D) := by
  induction t with
  | zero => simp [stock]
  | succ n ih => rw [gap_succ ρ r D S0 hρ, ih, pow_succ]; ring

/-- Cumulative virgin input equals `T·D·(1 - r)`. -/
theorem cumulative_eq (D r : ℝ) (T : ℕ) :
    cumulativeVirgin D r T = (T : ℝ) * D * (1 - r) := by
  simp [cumulativeVirgin, virgin, Finset.sum_const, Finset.card_range, nsmul_eq_mul]
  ring

/-- Convergence of the stock to the steady stock. -/
theorem stock_tendsto (ρ r D S0 : ℝ) (h0 : 0 ≤ ρ) (hρ : ρ < 1) :
    Tendsto (stock ρ r D S0) atTop (𝓝 (steadyStock ρ r D)) := by
  have hp : Tendsto (fun t : ℕ => ρ ^ t) atTop (𝓝 0) :=
    tendsto_pow_atTop_nhds_zero_of_lt_one h0 hρ
  have hg : Tendsto (fun t : ℕ => ρ ^ t * (S0 - steadyStock ρ r D) + steadyStock ρ r D)
      atTop (𝓝 (0 * (S0 - steadyStock ρ r D) + steadyStock ρ r D)) :=
    (hp.mul_const _).add_const _
  rw [zero_mul, zero_add] at hg
  refine hg.congr' (Eventually.of_forall fun t => ?_)
  have := gap_closed ρ r D S0 hρ t
  linarith

theorem main : ViridisStatements.P3.LoopMassBalance.Statement := by
  refine ⟨?_, ?_⟩
  · intro D r ρ S0 hD hr0 hr1 hρ0 hρ1
    refine ⟨⟨?_, ?_⟩, gap_closed ρ r D S0 hρ1, stock_tendsto ρ r D S0 hρ0 hρ1,
      cumulative_eq D r⟩
    · unfold virgin; exact mul_nonneg hD (by linarith)
    · unfold virgin; nlinarith
  · intro D T hD hT a _ b _ hab
    simp only [cumulative_eq]
    have hT' : (0 : ℝ) < T := by exact_mod_cast hT
    have : 0 < (T : ℝ) * D := mul_pos hT' hD
    nlinarith

end Candidate
\end{lstlisting}

\section{Results}
\label{sec:results}
We ran the program's checker on the candidate file. It runs the same stages as the verifier:
denylist, compile, statement match, axioms and kernel replay. The output was
\texttt{PASSED (preview)}, and the reported axioms were \texttt{propext},
\texttt{Classical.choice} and \texttt{Quot.sound}. The proof worked on the first attempt,
following the sketch recorded in the P3-T001 notebook. No automated prover was needed. The
certification page that precedes this paper gives the verifier's own judgement.

\section{What this settles and what it does not}
\label{sec:limits}
\paragraph{Settled.} Assuming the verifier accepts the file, the accounting identities of the
one-material loop model hold exactly as sealed. These are the virgin/recovered split, the fixed
point $S^\ast=rD/(1-\rho)$ with geometric convergence at rate $\rho$, the linear
cumulative-virgin identity, and strict monotonicity in $r$. A kernel built on these definitions
may rely on them for any parameter values in the stated ranges.

\paragraph{Not settled.}
\begin{itemize}
\item \textbf{No empirical content.} Nothing here estimates or constrains any real recovery
rate, retention rate, demand or material flow. Those remain field inputs. The theorem only says
what follows once they are given.
\item \textbf{Stylised dynamics.} Demand, $r$ and $\rho$ are constant. There is one material.
Recovery does not depend on the stock: the model does not require $rD\le\rho S(t)$, so recovered
supply is not limited by available stock. Quality loss, cascading and time lags are absent.
A stock-limited variant was flagged as an open question in P3-T001.
\item \textbf{Decoupling of $V$ and $S$.} In this model virgin input depends only on $r$, not on
the stock. The fixed point describes the size of the secondary stock, not the extraction
saving, which is $rD$ per period from the start.
\end{itemize}

\section{Reproducibility}
\sloppy
The proof file \texttt{proofs/P3.LoopMassBalance.lean} imports only \texttt{Mathlib} and the
sealed module \texttt{ViridisStatements.P3.LoopMassBalance}. It compiles under toolchain
\texttt{leanprover/lean4:v4.28.0} and the Mathlib commit below. The checker invocation is also
shown below. Alternatively, compile the file with \texttt{lean} and inspect
\texttt{\#print axioms Candidate.main} in a scratch file.
\begin{lstlisting}[language={},frame=none]
sealed module sha256:
  f00634594133f7edfb49f21bf30203d259cc8dfa7e62a40326bff1da1ea31ac7
Mathlib commit:
  8f9d9cff6bd728b17a24e163c9402775d9e6a365
viridis-lean-check proofs/P3.LoopMassBalance.lean P3.LoopMassBalance
\end{lstlisting}

\section{Changelog}
\begin{itemize}
\item 2026-11-05 (P3-T001): first version. Proposed P3.LoopMassBalance, P3.CascadedRecovery,
P3.RecoveryWorkFloor and P3.CyclicSymbiosis as statement modules, with controls, probe results
and rationale.
\item 2026-11-15 (P3-T002): P3.LoopMassBalance sealed (D-011). Narrowed the paper to a note on
this statement and added a full Lean proof and a human-readable proof. Moved the other three
proposals out of the paper body until they are sealed and proved. Dropped the two canon
references, which this note does not use.
\end{itemize}

\bibliographystyle{plainnat}
\bibliography{refs}
\end{document}
